\subsubsection{get\_llm}
\label{sec:get_llm}

\paragraph{Descrizione}
La funzione \texttt{get\_llm} recupera le infomrazioni relative a un LLM salvato.

\paragraph{Dettagli dell'endpoint}
\begin{itemize}
    \item \textbf{HTTP Method}: GET;
    \item \textbf{Endpoint}: /llms/\texttt{<id\_llm>};
    \item \textbf{Response Model}: LLMDTO;
    \item \textbf{Dependency Injection}:
    \begin{itemize}
        \item \textbf{llm\_service (LLMUseCase)}: la dipendenza viene risolta tramite \textit{Dependency Injector}. Questo permette di ottenere l'istanza del \texttt{LLMService} dal \texttt{Container}, che si occupa di tutte le configurazioni e inizializzazioni necessarie per il servizio.
    \end{itemize}
\end{itemize}

\paragraph{Parametri di input}
\begin{itemize}
    \item \textbf{id (int)}: ID dell'LLM salvato da ottenere.
\end{itemize}

\paragraph{Implementazione}
\begin{itemize}
    \item Restituisce lo stato dell'operazione e nel caso di successo un oggetto \texttt{LLMDTO} contenente i dati richiesti.
\end{itemize}


\subsubsection{get\_all\_llms}
\label{sec:get_all_llms}

\paragraph{Descrizione}
La funzione \texttt{get\_all\_llms} recupera la lista di tutti gli LLM salvati.

\paragraph{Dettagli dell'endpoint}
\begin{itemize}
    \item \textbf{HTTP Method}: GET;
    \item \textbf{Endpoint}: /llms;
    \item \textbf{Response Model}: LLMListDto;
    \item \textbf{Dependency Injection}:
    \begin{itemize}
        \item \textbf{llm\_service (LLMUseCase)}: la dipendenza viene risolta tramite \textit{Dependency Injector}. Questo permette di ottenere l'istanza del \texttt{LLMService} dal \texttt{Container}, che si occupa di tutte le configurazioni e inizializzazioni necessarie per il servizio.
    \end{itemize}
\end{itemize}

\paragraph{Implementazione}
\begin{itemize}
    \item Restituisce lo stato dell'operazione e nel caso di successo un oggetto \texttt{LLMListDTO} contenente i dati richiesti.
\end{itemize}



\subsubsection{delete\_llm}
\label{sec:delete_llm}

\paragraph{Descrizione}
La funzione \texttt{delete\_llm} elimina l'LLM salvato indicato.

\paragraph{Dettagli dell'endpoint}
\begin{itemize}
    \item \textbf{HTTP Method}: DELETE;
    \item \textbf{Endpoint}: /llms/\texttt{<id\_llm>};
    \item \textbf{Dependency Injection}:
    \begin{itemize}
        \item \textbf{llm\_service (LLMUseCase)}: la dipendenza viene risolta tramite \textit{Dependency Injector}. Questo permette di ottenere l'istanza del \texttt{LLMService} dal \texttt{Container}, che si occupa di tutte le configurazioni e inizializzazioni necessarie per il servizio.
    \end{itemize}
\end{itemize}

\paragraph{Parametri di input}
\begin{itemize}
    \item \textbf{id (int)}: ID dell'LLM salvato da eliminare.
\end{itemize}

\paragraph{Implementazione}
\begin{itemize}
    \item Restituisce lo stato dell'operazione.
\end{itemize}

\subsubsection{create\_llm}
\label{sec:create_llm}

\paragraph{Descrizione}
La funzione \texttt{create\_llm} salva un nuovo LLM.

\paragraph{Dettagli dell'endpoint}
\begin{itemize}
    \item \textbf{HTTP Method}: POST;
    \item \textbf{Endpoint}: \texttt{/llms};
    \item \textbf{Response Model}: LLMDto;
    \item \textbf{Dependency Injection}:
    \begin{itemize}
        \item \textbf{llm\_service (LLMUseCase)}: la dipendenza viene risolta tramite \textit{Dependency Injector}. Questo permette di ottenere l'istanza del \texttt{LLMService} dal \texttt{Container}, che si occupa di tutte le configurazioni e inizializzazioni necessarie per il servizio.
    \end{itemize}
\end{itemize}

\paragraph{Parametri di input}
\begin{itemize}
    \item \textbf{LLMDTO}: oggetto che rappresenta l'LLM da salvare.
\end{itemize}

\paragraph{Implementazione}
\begin{itemize}
    \item Restituisce lo stato dell'operazione e nel caso in cui sia andata a buon fine un oggetto \texttt{LLMDTO} che rappresenta LLM salvato.
\end{itemize}

\subsubsection{update\_llm}
\label{sec:update_LLM}

\paragraph{Descrizione}
La funzione \texttt{update\_llm} aggiorna un LLM salvato.

\paragraph{Dettagli dell'endpoint}
\begin{itemize}
    \item \textbf{HTTP Method}: PUT;
    \item \textbf{Endpoint}: \texttt{/llms/<id\_llm>};
    \item \textbf{Response Model}: LLMDTO;
    \item \textbf{Dependency Injection}:
    \begin{itemize}
        \item \textbf{llm\_service (LLMUseCase)}: la dipendenza viene risolta tramite \textit{Dependency Injector}. Questo permette di ottenere l'istanza del \texttt{LLMService} dal \texttt{Container}, che si occupa di tutte le configurazioni e inizializzazioni necessarie per il servizio.
    \end{itemize}
\end{itemize}

\paragraph{Parametri di input}
\begin{itemize}
    \item \textbf{id\_llm}: ID dell'LLM da aggiornare;
    \item \textbf{LLMDTO}: oggetto che rappresenta l'LLM da utilizzare nell'aggiornamento.
\end{itemize}

\paragraph{Implementazione}
\begin{itemize}
    \item Restituisce lo stato dell'operazione e nel caso in cui sia andata a buon fine un oggetto \texttt{LLMDTO} che rappresenta LLM aggiornato.
\end{itemize}

\paragraph{Tracciamento dei requisiti}
\begin{tabularx}{\textwidth}{|X|c|}
    \hline
    \textbf{ID} & \textbf{Funzione} \\
    \hline
        RFF-17.01, RFF-17.03, RFF-17.04, RFF-17.06, RFF-18.01 & get\_llm\\
   \hline
        RFF-16.01, RFF-16.05 & get\_all\_llms\\
   \hline
        RFF-16.04, RFF-17.02, RFF-17.05, RFF-17.07, RFF-17.08, RFF-18.02, RFF-18.04  & create\_llm\\
   \hline
        RFF-16.03, RFF-17.01, RFF-17.02, RFF-17.04, RFF-17.05, RFF-17.06, RFF-17.07, RFF-17.08, RFF-18.02, RFF-18.03, RFF-18.04 & update\_llm\\
   \hline
        RFF-16.02 & delete\_llm\\
   \hline
\end{tabularx}